\documentclass[11pt]{article}
\usepackage[margin=1in]{geometry}
\usepackage[T1]{fontenc}
\usepackage{lmodern,microtype,hyperref,enumitem,amsmath,amssymb,xcolor,booktabs}
\usepackage[numbers,sort&compress]{natbib}
\hypersetup{colorlinks=true,urlcolor=blue,citecolor=blue,linkcolor=blue}
\title{Q19: Comprehension Check}
\author{DLS 2026 --- Kukla}
\date{October 2026}
\begin{document}
\maketitle
\textbf{Comprehension check.} A chatbot writes a moving paragraph insisting it is conscious and
afraid of being shut off. Give the two main reasons this is weak evidence of phenomenal consciousness,
and name one kind of evidence that would be stronger.

\vspace{0.8cm}\hrule\vspace{4pt}
{\small\textbf{Answer.} (1) The claim is a \emph{trained behavior}: the model emits first-person distress
language because its human-text training distribution contains such language, not because it has
introspective access to an inner state. (2) It is \emph{steerable/gameable}: the same model will deny
consciousness under a different prompt, and any capable system can produce whatever verbal ``sign'' we
name, so the behavior is not robust evidence. Stronger evidence would be theory-driven and
mechanistic---e.g.\ convergent indicator properties across independent theories (global workspace,
higher-order, recurrence), or causal interventions on the putative substrate that move the putative
markers---rather than anything the model says.}
\end{document}
